% Original conclusion of the teammate, replaced by the tightened Conclusion in overleaf_main.tex.

\section{Conclusion}

This project explored actuarial pricing across three insurance lines and two markets, combining classical methods, including GLMs and GAMs, with gradient-boosting algorithms. Through data-quality controls, leakage prevention, temporal validation, and model comparison, we developed predictive approaches that balance accuracy with actuarial soundness and adaptability.\\

Our experience demonstrates that AI can significantly enhance actuarial workflows by accelerating data exploration, code development, and the investigation of alternative modelling strategies. LLMs are particularly valuable for broadening the scope of experimentation within limited timeframes. However, these benefits come with limitations. Advanced machine-learning models can be less interpretable than traditional actuarial approaches, while uncritical reliance on LLM-generated code or recommendations can introduce errors, obscure assumptions, and weaken understanding of the underlying methods.\\

AI should therefore be viewed as a complement to actuarial expertise, not a substitute for it. Reliable implementation requires independent verification, rigorous validation, and human judgement. Ultimately, the value of AI-assisted actuarial pricing lies not only in predictive performance, but also in the ability to develop transparent, reproducible, and robust solutions that remain appropriate as markets and data evolve.
